\vspace{0.25cm}
Dans la partie « Diffusion thermique», la mise en équations de la diffusion thermique est limitée au cas des solides ; on peut étendre les résultats ainsi établis aux milieux fluides en l'absence de convection en affirmant la généralisation des équations obtenues dans les solides. La loi phénoménologique de Newton à l'interface entre un solide et un fluide peut être utilisée dès lors qu'elle est fournie.
% \begin{tabularx}{\textwidth}{|X|X|}
% \hline Notions et contenus & Capacités exigibles \\
% \hline \multicolumn{2}{|l|}{ 3.3. Diffusion thermique } \\
% \hline Vecteur densité de flux thermique ja & Exprimer le flux thermique à travers une surface orientée en utilisant le vecteur jQ. \\
% \hline Premier principe de la thermodynamique. & Établir, pour un milieu solide, l'équation locale traduisant le premier principe dans le cas d'un problème ne dépendant qu'une d'une seule coordonnée d'espace en coordonnées cartésiennes, cylindriques et sphériques, éventuellement en présence de sources internes. Utiliser l'opérateur divergence et son expression fournie pour exprimer le bilan local dans le cas d'une géométrie quelconque, éventuellement en présence de sources internes. \\
% \hline Loi de Fourier. & Utiliser la Ioi de Fourier. Citer quelques ordres de grandeur de conductivité thermique dans les conditions usuelles : air, eau, béton, métaux. \\
% \hline
{\tiny
\begin{minipage}[t]{0.45\textwidth}
\begin{tabularx}{\linewidth}{|p{1.6cm}|X|}
\hline Notions et contenus & Capacités exigibles \\
\hline \multicolumn{2}{|l|}{ 3.3. Diffusion thermique } \\
\hline Vecteur densité de flux thermique $\Vec{\jth}$ & Exprimer le flux thermique à travers une surface orientée en utilisant le vecteur $\Vec{\jth}$. \\
\hline Premier principe de la thermodynamique. & Établir, pour un milieu solide, l'équation locale traduisant le premier principe dans le cas d'un problème ne dépendant qu'une d'une seule coordonnée d'espace en coordonnées cartésiennes, cylindriques et sphériques, éventuellement en présence de sources internes. Utiliser l'opérateur divergence et son expression fournie pour exprimer le bilan local dans le cas d'une géométrie quelconque, éventuellement en présence de sources internes. \\
\hline Loi de Fourier. & Utiliser la Ioi de Fourier. Citer quelques ordres de grandeur de conductivité thermique dans les conditions usuelles : air, eau, béton, métaux. \\
\hline
\end{tabularx}
\vfill
\end{minipage}%
\begin{minipage}[t]{0.55\textwidth}
\begin{tabularx}{\linewidth}{|p{1cm}|X|}
\hline Equation de la diffusion thermique. & Établir une équation de diffusion thermique. Utiliser l'opérateur laplacien et son expression fournie pour écrire l'équation de diffusion dans le cas d'une géométrie quelconque. Analyser une équation de diffusion en ordres de grandeur pour relier des échelles caractéristiques spatiale et temporelle. Utiliser la loi de Newton fournie comme condition aux limites à une interface solide-fluide. 


\udl{Capacité numérique} : à l'aide d'un langage de programmation, résoudre l'équation de la diffusion thermique à une dimension par une méthode des différences finies dérivée de la méthode d'Euler explicite de résolution des équations différentielles ordinaires. Mettre en œuvre un dispositif expérimental utilisant une caméra thermique ou un capteur dans le domaine des infrarouges. \\
\hline Régimes stationnaires. Résistance thermique. & Utiliser la conservation du flux thermique sous forme locale ou globale en l'absence de source interne. Définir la notion de résistance thermique par analogie avec l'électrocinétique. Établir l'expression d'une résistance thermique dans le cas d'un modèle unidimensionnel. Utiliser les lois d'associations de résistances thermiques. \\
\hline
\end{tabularx}
\vfill
\end{minipage}
}